% formal/dyntickrcu.tex
% mainfile: ../perfbook.tex
% SPDX-License-Identifier: CC-BY-SA-3.0

% Disable frame around VerbatimN in two-column layout
\IfTwoColumn{
\RecustomVerbatimEnvironment{VerbatimN}{Verbatim}%
{numbers=left,numbersep=5pt,xleftmargin=10pt,xrightmargin=0pt,frame=none}
\setlength{\lnlblraise}{0pt}
}{}

\subsection{Promela寓言：
			     {dynticks}与可抢占RCU}
\label{sec:formal:Promela Parable: dynticks and Preemptible RCU}

2008年初，一种可抢占的RCU变体被接受进入Linux主线内核，
以支持实时工作负载。
该变体类似于自2005年8月以来
\rt\ 补丁集~\cite{IngoMolnar05a}中的RCU实现。
实时工作负载需要可抢占RCU，因为早期的RCU实现在整个RCU
读侧临界区内禁用了抢占，导致过高的实时延迟。

然而，早期\rt\ 实现的一个缺点是每个\IX{grace period}
都需要在每个CPU上执行工作，即使该CPU处于低功耗的
``dynticks-idle''状态，因而无法执行RCU读侧临界区。
dynticks-idle状态背后的理念是，空闲的CPU应当被完全断电以
节省能源。
简而言之，可抢占RCU可能会使近期Linux内核中一项
有价值的节能特性失效。
尽管Josh Triplett和Paul McKenney曾讨论过一些方法，
允许CPU在整个RCU宽限期内保持低功耗状态
（从而保留Linux内核的节能能力），但直到Steve Rostedt
在\rt\ 补丁集中将一个新的dyntick实现与可抢占RCU集成后，
事情才到了非解决不可的地步。

这种组合导致Steve的一台系统在启动时挂起，因此在2007年十月，
Paul编写了一个对可抢占RCU宽限期处理的dynticks友好修改。
Steve编写了\co{rcu_irq_enter()}和\co{rcu_irq_exit()}
接口，从\co{irq_enter()}和\co{irq_exit()}中断
入口/出口函数中调用。
这些\co{rcu_irq_enter()}和\co{rcu_irq_exit()}
函数是必需的，以允许RCU可靠地处理dynticks-idle CPU
被瞬间唤醒以执行包含RCU读侧临界区的中断处理程序的情况。
有了这些修改，Steve的系统能够可靠启动。
但Paul假定我们不可能在第一次就把代码写对了，
因此继续定期检查代码。

Paul从2007年10月到2008年2月反复审查代码，
几乎每次都能发现至少一个bug。
有一次，Paul甚至在意识到一个bug是虚幻的之前，
就已经编写并测试了修复程序。
实际上在所有情况下，这些``bug''都被证明是虚幻的。

二月底时，Paul对这个游戏已经有点疲倦了。
因此他决定借助Promela和Spin的帮助。
以下内容呈现了一系列七个逐步增加复杂度的Promela模型，
最后一个通过了验证，状态空间消耗了约40\,GB的主内存。

更重要的是，Promela和Spin确实为我找到了一个非常微妙的bug！

\QuickQuiz{
	好吧，那真是太好了！
	现在，如果我碰巧没有一台配备40\,GB主内存的机器，
	我该怎么办？？？
}\QuickQuizAnswer{
	别急，这个问题有不少合法的答案：
	\begin{enumerate}
	\item	尝试编译器标志\co{-DCOLLAPSE}和\co{-DMA=N}
		以减少内存消耗。
		参见\cref{sec:formal:Running the QRCU Example}。
	\item	进一步优化模型，减少其内存消耗。
	\item	进行纸笔证明，也许可以从Linux内核代码中的
		注释开始。
	\item	设计细致的压力测试，虽然它们不能证明代码正确，
		但可以发现隐藏的bug。
	\item	目前有一些趋势是开发能在较小机器集群上进行
		模型检查的工具。
		但请注意，由于Paul偶尔能使用一些大型机器，
		我们实际上并没有使用过这类工具。
	\item	等待经济适用系统的内存容量扩展到能适应你
		的问题。
	\item	使用云计算服务之一来短期租用一台大型系统。
	\end{enumerate}
}\QuickQuizEnd

更好的做法是设计出一个更简单、更快速且状态空间更小的算法。
最好的情况是算法如此简单，其正确性对于一般观察者来说
显而易见！

\Crefrange{sec:formal:Introduction to Preemptible RCU and dynticks}
{sec:formal:Grace-Period Interface}
概述了可抢占RCU的dynticks接口，
随后
\cref{sec:formal:Validating Preemptible RCU and dynticks}
讨论了该接口的验证。

\subsubsection{可抢占RCU与dynticks简介}
\label{sec:formal:Introduction to Preemptible RCU and dynticks}

每CPU变量\co{dynticks_progress_counter}是dynticks与
可抢占RCU之间接口的核心。
当对应的CPU处于dynticks-idle模式时，该变量的值为偶数，
否则为奇数。
CPU退出dynticks-idle模式有以下三个原因：

\begin{enumerate}
\item	开始运行一个任务，
\item	进入可能嵌套的一组中断处理程序的最外层时，以及
\item	进入\IXacr{nmi}处理程序时。
\end{enumerate}

可抢占RCU的宽限期机制对\co{dynticks_progress_counter}
变量的值进行采样，以确定何时可以安全地忽略处于
dynticks-idle状态的CPU。

以下三节概述了任务接口、中断/NMI接口，
以及Linux内核v2.6.25-rc4中宽限期机制对
\co{dynticks_progress_counter}变量的使用。

\subsubsection{任务接口}
\label{sec:formal:Task Interface}

当某个CPU因为没有更多任务可运行而进入dynticks-idle模式时，
它会调用\co{rcu_enter_nohz()}：

\begin{VerbatimN}
static inline void rcu_enter_nohz(void)
{
	mb();
	__get_cpu_var(dynticks_progress_counter)++;
	WARN_ON(__get_cpu_var(dynticks_progress_counter) &
	        0x1);
}
\end{VerbatimN}

这个函数简单地递增\co{dynticks_progress_counter}并检查结果
是否为偶数，但在此之前先执行一个内存屏障，以确保任何其他
CPU在看到\co{dynticks_progress_counter}的新值时，也能看到
之前所有RCU读侧临界区的完成。

类似地，当处于dynticks-idle模式的CPU准备开始执行一个
新的可运行任务时，它会调用\co{rcu_exit_nohz()}：

\begin{VerbatimN}
static inline void rcu_exit_nohz(void)
{
	__get_cpu_var(dynticks_progress_counter)++;
	mb();
	WARN_ON(!(__get_cpu_var(dynticks_progress_counter) &
	          0x1));
}
\end{VerbatimN}

这个函数同样递增\co{dynticks_progress_counter}，
但在其后跟随一个内存屏障，以确保如果任何其他CPU
看到后续RCU读侧临界区的结果，那么该CPU也将看到
\co{dynticks_progress_counter}递增后的值。
最后，\co{rcu_exit_nohz()}检查递增的结果是否为奇数值。

\co{rcu_enter_nohz()}和\co{rcu_exit_nohz()}
函数处理CPU因任务执行而进入和退出dynticks-idle模式的情况，
但不处理中断，中断将在下一节中介绍。

\subsubsection{中断接口}
\label{sec:formal:Interrupt Interface}

\co{rcu_irq_enter()}和\co{rcu_irq_exit()}
函数分别处理中断/NMI的进入和退出。
当然，嵌套中断也必须被正确处理。
嵌套中断的可能性通过第二个每CPU变量\co{rcu_update_flag}
来处理，该变量在进入中断或NMI处理程序时（在
\co{rcu_irq_enter()}中）递增，
在退出时（在\co{rcu_irq_exit()}中）递减。
此外，已有的\co{in_interrupt()}原语用于区分最外层
中断和嵌套中断/NMI\@。

中断进入由如下所示的\co{rcu_irq_enter()}处理：

\begin{fcvlabel}[ln:formal:dyntickrcu:rcu_irq_enter]
\begin{VerbatimN}[commandchars=\\\[\]]
void rcu_irq_enter(void)
{
	int cpu = smp_processor_id();	\lnlbl[fetch]

	if (per_cpu(rcu_update_flag, cpu))	\lnlbl[inc:b]
		per_cpu(rcu_update_flag, cpu)++; \lnlbl[inc:e]
	if (!in_interrupt() &&			\lnlbl[chk_lv:b]
	    (per_cpu(dynticks_progress_counter,
	             cpu) & 0x1) == 0) {	\lnlbl[chk_lv:e]
		per_cpu(dynticks_progress_counter, cpu)++; \lnlbl[inc_cnt]
		smp_mb();			\lnlbl[mb]
		per_cpu(rcu_update_flag, cpu)++;\lnlbl[inc_flg]
	}
}
\end{VerbatimN}
\end{fcvlabel}

\begin{fcvref}[ln:formal:dyntickrcu:rcu_irq_enter]
\Clnref{fetch}获取当前CPU的编号，而\clnref{inc:b,inc:e}
在\co{rcu_update_flag}嵌套计数器已经非零时将其递增。
\Clnrefrange{chk_lv:b}{chk_lv:e}检查我们是否处于中断的
最外层，如果是，还检查\co{dynticks_progress_counter}
是否需要递增。
如果需要，\clnref{inc_cnt}递增\co{dynticks_progress_counter}，
\clnref{mb}执行内存屏障，\clnref{inc_flg}递增
\co{rcu_update_flag}。
与\co{rcu_exit_nohz()}一样，内存屏障确保任何其他CPU
在看到中断处理程序中RCU读侧临界区的效果时（在
\co{rcu_irq_enter()}调用之后），也能看到
\co{dynticks_progress_counter}的递增。
\end{fcvref}

\QuickQuizSeries{%
\QuickQuizB{
	为什么不简单地递增\co{rcu_update_flag}，然后仅在
	\co{rcu_update_flag}的旧值为零时才递增
	\co{dynticks_progress_counter}？？？
}\QuickQuizAnswerB{
	这在存在NMI的情况下会失败。
	要理解这一点，假设在\co{rcu_irq_enter()}递增了
	\co{rcu_update_flag}之后、但在递增
	\co{dynticks_progress_counter}之前收到了一个NMI。
	由NMI调用的\co{rcu_irq_enter()}实例将看到
	\co{rcu_update_flag}的原始值为非零，因此会放弃递增
	\co{dynticks_progress_counter}。
	这将使RCU宽限期机制无法知道NMI处理程序正在该CPU上
	执行，导致NMI处理程序中的任何RCU读侧临界区失去RCU保护。

	NMI处理程序的存在——按定义不可被屏蔽——确实使这段
	代码变得复杂。
}\QuickQuizEndB
%
\QuickQuizE{
	\begin{fcvref}[ln:formal:dyntickrcu:rcu_irq_enter]
	但是如果\clnref{chk_lv:b}发现我们处于最外层中断，
	我们不是\emph{总是}需要递增
	\co{dynticks_progress_counter}吗？
	\end{fcvref}
}\QuickQuizAnswerE{
	如果我们中断的是一个正在运行的任务则不需要！
	在那种情况下，\co{dynticks_progress_counter}已经被
	\co{rcu_exit_nohz()}递增过了，不需要再次递增。
}\QuickQuizEndE
}

中断退出由\co{rcu_irq_exit()}以类似方式处理：

\begin{fcvlabel}[ln:formal:dyntickrcu:rcu_irq_exit]
\begin{VerbatimN}[commandchars=\\\[\]]
void rcu_irq_exit(void)
{
	int cpu = smp_processor_id();	\lnlbl[fetch]

	if (per_cpu(rcu_update_flag, cpu)) { \lnlbl[chk_flg]
		if (--per_cpu(rcu_update_flag, cpu)) \lnlbl[dec_flg]
			return;
		WARN_ON(in_interrupt());	\lnlbl[verify]
		smp_mb();			\lnlbl[mb]
		per_cpu(dynticks_progress_counter, cpu)++; \lnlbl[inc_cnt]
		WARN_ON(per_cpu(dynticks_progress_counter, \lnlbl[vrf_even:b]
		                cpu) & 0x1); \lnlbl[vrf_even:e]
	}
}
\end{VerbatimN}
\end{fcvlabel}

\begin{fcvref}[ln:formal:dyntickrcu:rcu_irq_exit]
\Clnref{fetch}与之前一样获取当前CPU的编号。
\Clnref{chk_flg}检查\co{rcu_update_flag}是否非零，
如果为零则立即返回（通过从函数末尾落出）。
否则，\clnrefthro{dec_flg}{vrf_even:e}开始发挥作用。
\Clnref{dec_flg}递减\co{rcu_update_flag}，如果结果不为零
则返回。
\Clnref{verify}验证我们确实正在离开嵌套中断的最外层，
\clnref{mb}执行内存屏障，
\clnref{inc_cnt}递增\co{dynticks_progress_counter}，
\clnref{vrf_even:b,vrf_even:e}验证该变量现在是偶数。
与\co{rcu_enter_nohz()}一样，内存屏障确保任何其他CPU
在看到\co{dynticks_progress_counter}的递增时，
也能看到中断处理程序中RCU读侧临界区的效果
（在\co{rcu_irq_exit()}调用之前）。
\end{fcvref}

以上两节描述了在任务、中断和NMI进入和退出dynticks-idle
模式期间，\co{dynticks_progress_counter}变量是如何维护的。
下一节描述可抢占RCU的宽限期机制如何使用该变量。

\subsubsection{宽限期接口}
\label{sec:formal:Grace-Period Interface}

在\cref{fig:formal:Preemptible RCU State Machine}所示的
四个可抢占RCU宽限期状态中，
只有\co{rcu_try_flip_waitack_state}
和\co{rcu_try_flip_waitmb_state}状态需要等待
其他CPU响应。

\begin{figure}
\centering
\resizebox{2.5in}{!}{\includegraphics{formal/RCUpreemptStates}}
\caption{可抢占RCU状态机}
\label{fig:formal:Preemptible RCU State Machine}
\end{figure}

当然，如果给定的CPU处于dynticks-idle状态，我们不应该
等待它。
因此，在进入这两个状态之一之前，前一个状态会对每个CPU的
\co{dynticks_progress_counter}变量进行快照，并将快照
存放在另一个每CPU变量\co{rcu_dyntick_snapshot}中。
这通过调用如下所示的
\co{dyntick_save_progress_counter()}来完成：

\begin{VerbatimN}
static void dyntick_save_progress_counter(int cpu)
{
	per_cpu(rcu_dyntick_snapshot, cpu) =
		per_cpu(dynticks_progress_counter, cpu);
}
\end{VerbatimN}

\co{rcu_try_flip_waitack_state}状态调用如下所示的
\co{rcu_try_flip_waitack_needed()}：

\begin{fcvlabel}[ln:formal:dyntickrcu:rcu_try_flip_waitack_needed]
\begin{VerbatimN}[samepage=true,commandchars=\\\[\]]
static inline int
rcu_try_flip_waitack_needed(int cpu)
{
	long curr;
	long snap;

	curr = per_cpu(dynticks_progress_counter, cpu); \lnlbl[curr]
	snap = per_cpu(rcu_dyntick_snapshot, cpu); \lnlbl[snap]
	smp_mb();				\lnlbl[mb]
	if ((curr == snap) && ((curr & 0x1) == 0)) \lnlbl[chk_remain]
		return 0;			\lnlbl[ret_0_r]
	if ((curr - snap) > 2 || (snap & 0x1) == 0) \lnlbl[chk_idle]
		return 0;			\lnlbl[ret_0_i]
	return 1;				\lnlbl[ret_1]
}
\end{VerbatimN}
\end{fcvlabel}

\begin{fcvref}[ln:formal:dyntickrcu:rcu_try_flip_waitack_needed]
\Clnref{curr,snap}分别获取\co{dynticks_progress_counter}的
当前版本和快照版本。
\clnref{mb}上的内存屏障确保后续
\co{rcu_try_flip_waitzero_state}中的计数器检查
在这些计数器获取之后执行。
\Clnref{chk_remain,ret_0_r}在该CPU自快照获取以来一直
保持在dynticks-idle状态时返回零（意味着不需要与
指定CPU通信）。
类似地，\clnref{chk_idle,ret_0_i}在该CPU最初处于
dynticks-idle状态或已经完整地经历了一次
dynticks-idle状态时返回零。
在这两种情况下，CPU不可能保留宽限期计数器的旧值。
如果这两个条件都不满足，\clnref{ret_1}返回一，
意味着该CPU需要显式响应。
\end{fcvref}

对于\co{rcu_try_flip_waitmb_state}状态，它调用如下所示的
\co{rcu_try_flip_waitmb_needed()}：

\begin{fcvlabel}[ln:formal:dyntickrcu:rcu_try_flip_waitmb_needed]
\begin{VerbatimN}[commandchars=\\\[\]]
static inline int
rcu_try_flip_waitmb_needed(int cpu)
{
	long curr;
	long snap;

	curr = per_cpu(dynticks_progress_counter, cpu);
	snap = per_cpu(rcu_dyntick_snapshot, cpu);
	smp_mb();
	if ((curr == snap) && ((curr & 0x1) == 0))
		return 0;
	if (curr != snap)		\lnlbl[chk_to_from]
		return 0;		\lnlbl[ret_0]
	return 1;
}
\end{VerbatimN}
\end{fcvlabel}

\begin{fcvref}[ln:formal:dyntickrcu:rcu_try_flip_waitmb_needed]
这与\co{rcu_try_flip_waitack_needed()}非常相似，
不同之处在于\clnref{chk_to_from,ret_0}，因为任何进入或
离开dynticks-idle状态的转换都会执行
\co{rcu_try_flip_waitmb_state}状态所需的内存屏障。
\end{fcvref}

我们现在已经看到了RCU与dynticks-idle状态之间接口涉及的
所有代码。
下一节将构建用于验证这些代码的Promela模型。

\QuickQuiz{
	你能在本节的代码中发现任何bug吗？
}\QuickQuizAnswer{
	阅读下一节来看看你是否正确。
}\QuickQuizEnd

\subsection{验证可抢占RCU和dynticks}
\label{sec:formal:Validating Preemptible RCU and dynticks}

本节逐步开发dynticks与RCU之间接口的Promela模型，
\crefrange{sec:formal:Basic Model}{sec:formal:Validating NMI Handlers}
中的每一节都展示了一个步骤，从进程级代码开始，
逐步添加断言、中断，最后是NMI。

\Cref{sec:formal:Lessons (Re)Learned}列出了
在这一过程中（重新）学到的经验教训，
\crefrange{sec:formal:Simplicity Avoids Formal Verification}{sec:formal:Discussion}
则展示了一个更简单的RCU dynticks问题解决方案。

\subsubsection{基本模型}
\label{sec:formal:Basic Model}

本节将进程级的dynticks进入/退出代码和宽限期处理翻译为
Promela~\cite{Holzmann03a}。
我们从2.6.25-rc4内核的\co{rcu_exit_nohz()}和
\co{rcu_enter_nohz()}开始，
将它们放入一个Promela进程中，以循环方式建模退出和
进入dynticks-idle模式，如下所示：

\input{CodeSamples/formal/promela/dyntick/dyntickRCU-base@dyntick_nohz.fcv}

\begin{fcvref}[ln:formal:promela:dyntick:dyntickRCU-base:dyntick_nohz]
\Clnref{do,od}定义了一个循环。
\Clnref{break}在循环计数器~\co{i}超过限制
\co{MAX_DYNTICK_LOOP_NOHZ}时退出循环。
\Clnref{kick_loop}告诉循环构造在每次循环迭代中执行
\clnrefrange{ex_inc:b}{inc_i}。
因为\clnref{break,kick_loop}上的条件是互斥的，
Promela通常的随机选择真条件的行为被禁用。
\Clnref{ex_inc:b,ex_inc:e}建模了\co{rcu_exit_nohz()}对
\co{dynticks_progress_counter}的非原子递增，而
\clnref{ex_assert}建模了\co{WARN_ON()}。
\co{atomic}构造只是为了减少Promela的状态空间，
因为\co{WARN_ON()}严格来说不是算法的一部分。
\Clnrefrange{ent_inc:b}{ent_inc:e}类似地建模了
\co{rcu_enter_nohz()}的递增和\co{WARN_ON()}。
最后，\clnref{inc_i}递增循环计数器。
\end{fcvref}

因此，循环的每次迭代建模了一个CPU退出dynticks-idle模式
（例如，开始执行一个任务），然后重新进入dynticks-idle模式
（例如，该任务阻塞）。

\QuickQuizSeries{%
\QuickQuizB{
	为什么在Promela中不对\co{rcu_exit_nohz()}和
	\co{rcu_enter_nohz()}中的内存屏障进行建模？
}\QuickQuizAnswerB{
	Promela假设顺序一致性，因此没有必要对内存屏障建模。
	实际上，人们反而必须显式地对缺少内存屏障的情况建模，
	例如，如
	\cref{lst:formal:QRCU Unordered Summation}
	（\cpageref{lst:formal:QRCU Unordered Summation}）
	所示。
}\QuickQuizEndB
%
\QuickQuizE{
	先建模\co{rcu_exit_nohz()}再建模
	\co{rcu_enter_nohz()}是不是有点奇怪？
	先建模进入再建模退出不是更自然吗？
}\QuickQuizAnswerE{
	那样可能确实更自然，但我们在后面要添加的活性检查中
	需要这种特定的顺序。
}\QuickQuizEndE
}

下一步是建模RCU宽限期处理的接口。
为此，我们需要建模
\co{dyntick_save_progress_counter()}、
\co{rcu_try_flip_waitack_needed()}、
\co{rcu_try_flip_waitmb_needed()}，
以及\co{rcu_try_flip_waitack()}和
\co{rcu_try_flip_waitmb()}的部分内容，
均来自2.6.25-rc4内核。
下面的\co{grace_period()} Promela进程建模了这些函数
在可抢占RCU宽限期处理的单次遍历中的调用方式。

\input{CodeSamples/formal/promela/dyntick/dyntickRCU-base@grace_period.fcv}

\begin{fcvref}[ln:formal:promela:dyntick:dyntickRCU-base:grace_period]
\Clnrefrange{print:b}{print:e}打印循环限制
（但只在出错时打印到``.trail''文件中），
并建模了\co{rcu_try_flip_idle()}中的一行代码及其对
\co{dyntick_save_progress_counter()}的调用，
该调用对当前CPU的\co{dynticks_progress_counter}变量
进行快照。
这两行原子地执行以减少状态空间。

\Clnrefrange{do1}{od1}建模了\co{rcu_try_flip_waitack()}
中的相关代码及其对
\co{rcu_try_flip_waitack_needed()}的调用。
该循环建模了宽限期状态机等待每个CPU的计数器翻转确认，
但仅涉及与dynticks-idle CPU交互的部分。

\Clnref{snap}建模了\co{rcu_try_flip_waitzero()}中的
一行及其对\co{dyntick_save_progress_counter()}的调用，
再次对CPU的\co{dynticks_progress_counter}变量进行快照。

最后，\clnrefrange{do2}{od2}建模了
\co{rcu_try_flip_waitack()}中的相关代码及其对
\co{rcu_try_flip_waitack_needed()}的调用。
该循环建模了宽限期状态机等待每个CPU执行内存屏障，
但同样仅涉及与dynticks-idle CPU交互的部分。
\end{fcvref}

\QuickQuiz{
	等一下！
	在Linux内核中，\co{dynticks_progress_counter}和
	\co{rcu_dyntick_snapshot}都是每CPU变量。
	那为什么它们在这里被建模为单个全局变量？
}\QuickQuizAnswer{
	因为宽限期代码分别处理每个CPU的
	\co{dynticks_progress_counter}和
	\co{rcu_dyntick_snapshot}变量，
	我们可以将状态压缩到单个CPU上。
	如果宽限期代码在特定CPU上针对特定值做特殊处理，
	那么我们确实需要建模多个CPU。
	但幸运的是，我们可以安全地限制在两个CPU上：
	一个运行宽限期处理，另一个进入和离开dynticks-idle模式。
}\QuickQuizEnd

得到的模型（\path{dyntickRCU-base.spin}），
使用\path{runspin.sh}脚本运行时，
生成了691个状态并且无错误通过，
鉴于该模型完全缺少能够发现故障的断言，
这一点丝毫不令人意外。
因此，下一节添加安全性断言。

\subsubsection{验证安全性}
\label{sec:formal:Validating Safety}

一个安全的RCU实现绝不能允许宽限期在任何在宽限期开始
之前启动的RCU读者完成之前结束。
这通过一个\co{grace_period_state}变量来建模，
该变量可以取以下三种状态：

\input{CodeSamples/formal/promela/dyntick/dyntickRCU-base-s@grace_period_state.fcv}

\co{grace_period()}进程在经过宽限期各阶段时设置该变量，
如下所示：

\input{CodeSamples/formal/promela/dyntick/dyntickRCU-base-s@grace_period.fcv}

\begin{fcvref}[ln:formal:promela:dyntick:dyntickRCU-base-s:grace_period]
\Clnref{upd_gps1,upd_gps2,upd_gps3,upd_gps4,upd_gps5,upd_gps6}
更新该变量（在可行的地方与算法操作原子地组合），
以允许\co{dyntick_nohz()}进程验证基本的RCU安全性属性。
此验证的形式是断言\co{grace_period_state}变量的值
不能在RCU读者可能持续存在的时间段内从
\co{GP_IDLE}跳到\co{GP_DONE}。
\end{fcvref}

\QuickQuiz{
	\begin{fcvref}[ln:formal:promela:dyntick:dyntickRCU-base-s:grace_period]
	既然在\clnref{upd_gps3,upd_gps4}上有一对背靠背的
	\co{grace_period_state}修改，
	我们如何确保\clnref{upd_gps3}的修改不会丢失？
	\end{fcvref}
}\QuickQuizAnswer{
	回想一下Promela和Spin会追踪每一种可能的状态变化序列。
	因此，时序是无关紧要的：
	Promela/Spin会很乐意在这两条语句之间插入模型的
	所有剩余部分，除非某个状态变量明确禁止这样做。
}\QuickQuizEnd

\co{dyntick_nohz()} Promela进程实现此验证如下所示：

\input{CodeSamples/formal/promela/dyntick/dyntickRCU-base-s@dyntick_nohz.fcv}

\begin{fcvref}[ln:formal:promela:dyntick:dyntickRCU-base-s:dyntick_nohz]
\Clnref{new_flg}在任务执行开始时如果
\co{grace_period_state}变量的值为\co{GP_IDLE}，
则设置一个新的\co{old_gp_idle}标志，
而\clnref{assert:b,assert:e}处的断言在
\co{grace_period_state}变量在任务执行期间
推进到\co{GP_DONE}时触发，
这是非法的，因为单个RCU读侧临界区可能跨越整个
中间时间段。
\end{fcvref}

得到的模型（\path{dyntickRCU-base-s.spin}），
使用\path{runspin.sh}脚本运行时，
生成了964个状态并无错误通过，这令人放心。
尽管如此，虽然安全性至关重要，避免宽限期无限停滞也同样
重要。因此下一节将介绍活性验证。

\subsubsection{验证活性}
\label{sec:formal:Validating Liveness}

\begin{fcvref}[ln:formal:promela:dyntick:dyntickRCU-base-sl-busted:dyntick_nohz]
虽然活性可能很难证明，但这里有一个简单的技巧适用。
第一步是让\co{dyntick_nohz()}通过一个
\co{dyntick_nohz_done}变量来指示它已完成，
如下面的\clnref{done}所示：
\end{fcvref}

\input{CodeSamples/formal/promela/dyntick/dyntickRCU-base-sl-busted@dyntick_nohz.fcv}

有了这个变量，我们可以在\co{grace_period()}中添加断言
来检查是否存在不必要的阻塞，如下所示：

\input{CodeSamples/formal/promela/dyntick/dyntickRCU-base-sl-busted@grace_period.fcv}

\begin{fcvref}[ln:formal:promela:dyntick:dyntickRCU-base-sl-busted:grace_period]
我们在\clnref{shex}添加了\co{shouldexit}变量，
在\clnref{init_shex}将其初始化为零。
\Clnref{assert_shex}断言\co{shouldexit}未被设置，
而\clnref{set_shex}将\co{shouldexit}设置为
\co{dyntick_nohz()}维护的\co{dyntick_nohz_done}变量。
因此，如果我们在\co{dyntick_nohz()}执行完成后
尝试在等待计数器翻转确认的循环中再次迭代，
该断言就会触发。
毕竟，如果\co{dyntick_nohz()}已完成，就不会再有状态
变化来强制我们退出循环，因此在这种状态下迭代两次
意味着无限循环，进而意味着宽限期永远不会结束。

\Clnref{init_shex2,assert_shex2,set_shex2}对第二个
（内存屏障）循环以类似方式操作。

然而，运行该模型（\path{dyntickRCU-base-sl-busted.spin}）
会导致失败，因为\clnref{chk_2}检查了错误的变量
是否为偶数。
失败时，\co{spin}会写出一个``trail''文件
（\path{dyntickRCU-base-sl-busted.spin.trail}），
记录导致失败的状态序列。
使用``\co{spin -t -p -g -l\ }%
\path{dyntickRCU-base-sl-busted.spin}''
命令可以让\co{spin}重新追踪这个状态序列，
打印执行的语句和变量的值
（\path{dyntickRCU-base-sl-busted.spin.trail.txt}）。
请注意，由于\co{spin}将两个函数放在单个文件中，
行号与上面的清单不匹配。
但是，行号\emph{确实}与完整模型
（\path{dyntickRCU-base-sl-busted.spin}）匹配。

我们看到\co{dyntick_nohz()}进程在步骤34完成
（搜索``34:''），但\co{grace_period()}进程仍然
未能退出循环。
\co{curr}的值为\co{6}（见步骤35），
\co{snap}的值为\co{5}（见步骤17）。
因此\clnref{chk_1}上的第一个条件不成立，因为
\qco{curr != snap}，而\clnref{chk_2}上的第二个条件
也不成立，因为\co{snap}是奇数且\co{curr}仅比
\co{snap}大一。
\end{fcvref}

所以这两个条件中必有一个是不正确的。
参考\co{rcu_try_flip_waitack_needed()}中
第一个条件的注释块：

\begin{quote}
	如果CPU在整个期间保持在dynticks模式且没有接收任何
	中断、NMI、SMI或其他，那么它不可能处于
	\co{rcu_read_lock()}的中间，所以它执行的下一个
	\co{rcu_read_lock()}必须使用计数器的新值。
	因此我们可以安全地假装该CPU已经确认了计数器。
\end{quote}

第一个条件确实与此匹配，因为如果\qco{curr == snap}
且\co{curr}为偶数，那么对应的CPU在整个期间一直处于
dynticks-idle模式，符合要求。
那么我们来看看第二个条件的注释块：

\begin{quote}
	如果CPU在没有活跃irq处理程序的情况下经历了或进入了
	dynticks空闲阶段，那么如上所述，我们可以安全地假装
	该CPU已经确认了计数器。
\end{quote}

条件的第一部分是正确的，因为如果\co{curr}和\co{snap}
相差二，则中间至少有一个偶数，对应于完整地经历了一个
dynticks-idle阶段。
然而，条件的第二部分对应的是\emph{开始}于dynticks-idle
模式，而不是\emph{结束}于该模式。
因此我们需要测试的是\co{curr}而非\co{snap}是否为偶数。

修正后的C代码如下：

\begin{fcvlabel}[ln:formal:dyntickrcu:rcu_try_flip_waitack_needed_fixed]
\begin{VerbatimN}[commandchars=\\\[\]]
static inline int
rcu_try_flip_waitack_needed(int cpu)
{
	long curr;
	long snap;

	curr = per_cpu(dynticks_progress_counter, cpu);
	snap = per_cpu(rcu_dyntick_snapshot, cpu);
	smp_mb();
	if ((curr == snap) && ((curr & 0x1) == 0)) \lnlbl[if:b]
		return 0;
	if ((curr - snap) > 2 || (curr & 0x1) == 0)
		return 0;			\lnlbl[if:e]
	return 1;
}
\end{VerbatimN}
\end{fcvlabel}

\begin{fcvref}[ln:formal:dyntickrcu:rcu_try_flip_waitack_needed_fixed]
\Clnrefrange{if:b}{if:e}现在可以合并和简化，
得到如下结果。
类似的简化也可以应用于\co{rcu_try_flip_waitmb_needed()}。
\end{fcvref}

\begin{VerbatimN}[commandchars=\\\[\]]
static inline int
rcu_try_flip_waitack_needed(int cpu)
{
	long curr;
	long snap;

	curr = per_cpu(dynticks_progress_counter, cpu);
	snap = per_cpu(rcu_dyntick_snapshot, cpu);
	smp_mb();
	if ((curr - snap) >= 2 || (curr & 0x1) == 0)
		return 0;
	return 1;
}
\end{VerbatimN}

在模型（\path{dyntickRCU-base-sl.spin}）中进行相应修正后，
得到一个包含661个状态的正确验证，无错误通过。
然而，值得注意的是，活性验证的第一个版本未能捕获这个bug，
原因是活性验证本身存在一个bug。
这个活性验证的bug是通过在\co{grace_period()}进程中
插入一个无限循环，并注意到活性验证代码未能检测到
该问题而发现的！

我们现在已经成功验证了安全性和活性条件，但仅限于进程运行
和阻塞的情况。
我们还需要处理中断，这是下一节要讨论的任务。

\subsubsection{中断}
\label{sec:formal:Interrupts}

在Promela中有两种方式来建模中断：
\begin{enumerate}
\item	使用C预处理器技巧在\co{dynticks_nohz()}进程的
	每对相邻语句之间插入中断处理程序，或者
\item	用一个单独的进程来建模中断处理程序。
\end{enumerate}

经过一番思考，第二种方法的状态空间会更小，
尽管它要求中断处理程序相对于\co{dynticks_nohz()}进程
原子执行，但不需要相对于\co{grace_period()}进程
原子执行。

幸运的是，Promela允许从原子语句中分支跳出。
这个技巧允许我们让中断处理程序设置一个标志，
并重新编写\co{dynticks_nohz()}以原子地检查该标志，
仅在标志未设置时才执行。
这可以通过一个接受标签和Promela语句的C预处理器宏来实现，
如下所示：

\input{CodeSamples/formal/promela/dyntick/dyntickRCU-irqnn-ssl@EXECUTE_MAINLINE.fcv}

可以如下使用该宏：

\begin{VerbatimU}
EXECUTE_MAINLINE(stmt1,
                 tmp = dynticks_progress_counter)
\end{VerbatimU}

\begin{fcvref}[ln:formal:promela:dyntick:dyntickRCU-irqnn-ssl:EXECUTE_MAINLINE]
宏的\clnref{label}创建指定的语句标签。
\Clnrefrange{atm:b}{atm:e}是一个原子块，测试
\co{in_dyntick_irq}变量，如果该变量被设置
（表示中断处理程序处于活动状态），则从原子块跳出
回到标签处。
否则，\clnref{else}执行指定的语句。
总体效果是，主线执行在中断活动时暂停，这正是所需要的。
\end{fcvref}

\subsubsection{验证中断处理程序}
\label{sec:formal:Validating Interrupt Handlers}

第一步是将\co{dyntick_nohz()}转换为
\co{EXECUTE_MAINLINE()}形式，如下所示：

\input{CodeSamples/formal/promela/dyntick/dyntickRCU-irqnn-ssl@dyntick_nohz.fcv}

\begin{fcvref}[ln:formal:promela:dyntick:dyntickRCU-irqnn-ssl:dyntick_nohz]
需要注意的是，当一组语句被传递给\co{EXECUTE_MAINLINE()}时，
如\clnrefrange{stmt2:b}{stmt2:e}，该组中的所有语句都原子地执行。
\end{fcvref}

\QuickQuizSeries{%
\QuickQuizB{
	但是如果你需要单个\co{EXECUTE_MAINLINE()}组中的
	语句非原子地执行，你会怎么做？
}\QuickQuizAnswerB{
	最简单的做法是将每个这样的语句放入自己的
	\co{EXECUTE_MAINLINE()}语句中。
}\QuickQuizEndB
%
\QuickQuizE{
	但是如果\co{dynticks_nohz()}进程有带条件的
	\qco{if}或\qco{do}语句，而这些构造的语句体
	需要非原子地执行呢？
}\QuickQuizAnswerE{
	一种方法，正如我们将在后面的章节中看到的，
	是使用显式标签和\qco{goto}语句。
	例如，以下构造：

\begin{VerbatimU}
	if
	:: i == 0 -> a = -1;
	:: else -> a = -2;
	fi;
\end{VerbatimU}
%
	可以建模为类似如下的形式：

\begin{VerbatimU}
	EXECUTE_MAINLINE(stmt1,
	                 if
	                 :: i == 0 -> goto stmt1_then;
	                 :: else -> goto stmt1_else;
	                 fi)
	stmt1_then: skip;
	EXECUTE_MAINLINE(stmt1_then1, a = -1; goto stmt1_end)
	stmt1_else: skip;
	EXECUTE_MAINLINE(stmt1_then1, a = -2)
	stmt1_end: skip;
\end{VerbatimU}

	然而，在\qco{if}语句的情况下，该宏的帮助并不明显，
	因此这类情况将在以下各节中直接展开编写。
}\QuickQuizEndE
}

下一步是编写\co{dyntick_irq()}进程来建模中断处理程序：

\input{CodeSamples/formal/promela/dyntick/dyntickRCU-irqnn-ssl@dyntick_irq.fcv}

\begin{fcvref}[ln:formal:promela:dyntick:dyntickRCU-irqnn-ssl:dyntick_irq]
\clnrefrange{do}{od}的循环建模了最多\co{MAX_DYNTICK_LOOP_IRQ}
次中断，\clnref{cond1,cond2}构成循环条件，
\clnref{inc_i}递增控制变量。
\Clnref{in_irq}告诉\co{dyntick_nohz()}中断处理程序正在
运行，\clnref{clr_in_irq}告诉\co{dyntick_nohz()}该处理
程序已完成。
\Clnref{irq_done}用于活性验证，与\co{dyntick_nohz()}中
对应的行相同。
\end{fcvref}

\QuickQuiz{
	\begin{fcvref}[ln:formal:promela:dyntick:dyntickRCU-irqnn-ssl:dyntick_irq]
	为什么\clnref{clr_in_irq,inc_i}（\qtco{in_dyntick_irq = 0;}
	和\qco{i++;}）是原子执行的？
	\end{fcvref}
}\QuickQuizAnswer{
	这些代码行属于控制模型，而不是被建模的代码，
	因此没有理由对它们进行非原子建模。
	将它们原子化建模的动机是减少状态空间的大小。
}\QuickQuizEnd

\begin{fcvref}[ln:formal:promela:dyntick:dyntickRCU-irqnn-ssl:dyntick_irq]
\Clnrefrange{enter:b}{enter:e}建模\co{rcu_irq_enter()}，
\clnref{add_prmt_cnt:b,add_prmt_cnt:e}建模\co{__irq_enter()}
的相关片段。
\Clnrefrange{vrf_safe:b}{vrf_safe:e}以与\co{dynticks_nohz()}
中对应行相同的方式验证安全性。
\Clnref{irq_exit:b,irq_exit:e}建模\co{__irq_exit()}的
相关片段，最后\clnrefrange{exit:b}{exit:e}建模
\co{rcu_irq_exit()}。
\end{fcvref}

\QuickQuiz{
	这个\co{dynticks_irq()}进程无法建模中断的
	哪个属性？
}\QuickQuizAnswer{
	一个这样的属性是嵌套中断，
	下一节将处理这个问题。
}\QuickQuizEnd

\co{grace_period()}进程变为如下所示：

\input{CodeSamples/formal/promela/dyntick/dyntickRCU-irqnn-ssl@grace_period.fcv}

\begin{fcvref}[ln:formal:promela:dyntick:dyntickRCU-irqnn-ssl:grace_period]
\co{grace_period()} 的实现与之前的非常相似。
唯一的变化是在 \clnref{MDLI} 上添加了新的中断计数参数，
在 \clnref{edit1,edit3} 上将新的 \co{dyntick_irq_done}
变量添加到活性检查中，当然还有 \clnref{edit2,edit4} 上的优化。
\end{fcvref}

此模型（\path{dyntickRCU-irqnn-ssl.spin}）
以大约五十万个状态产生了正确的验证，无错误通过。
然而，此版本的模型不处理嵌套中断。
这个主题将在下一节中讨论。

\subsubsection{验证嵌套中断处理程序}
\label{sec:formal:Validating Nested Interrupt Handlers}

嵌套中断处理程序可以通过拆分 \co{dyntick_irq()} 中循环的
主体来建模，如下所示：

\input{CodeSamples/formal/promela/dyntick/dyntickRCU-irq-ssl@dyntick_irq.fcv}

\begin{fcvref}[ln:formal:promela:dyntick:dyntickRCU-irq-ssl:dyntick_irq]
这与之前的 \co{dynticks_irq()} 进程类似。
它在 \clnref{j} 上添加了第二个计数器变量 \co{j}，使得
\co{i} 计数中断处理程序的进入次数，\co{j} 计数退出次数。
\clnref{om} 上的 \co{outermost} 变量帮助确定何时需要
为安全检查采样 \co{grace_period_state} 变量。
\clnref{chk_ex:b,chk_ex:e} 上的循环退出检查被更新为要求
指定数量的中断处理程序既被进入也被退出，
\co{i} 的递增被移到 \clnref{inc_i}，即中断进入模型的末尾。
\Clnrefrange{atm1:b}{atm1:e} 设置 \co{outermost} 变量以指示
这是否是一组嵌套中断中最外层的，
并设置 \co{dyntick_nohz()} 进程使用的
\co{in_dyntick_irq} 变量。
\Clnrefrange{atm2:b}{atm2:e} 捕获 \co{grace_period_state}
变量的状态，但仅在最外层中断处理程序中执行。

\Clnref{cnd_ex} 有中断退出建模的 do 循环条件：
只要我们退出的中断比进入的少，就可以退出另一个中断。
\Clnrefrange{atm3:b}{atm3:e}
检查安全准则，但仅在从最外层中断级别退出时执行。
最后，\clnrefrange{atm4:b}{atm4:e} 递增中断退出计数 \co{j}，
如果这是最外层中断级别，则清除 \co{in_dyntick_irq}。
\end{fcvref}

此模型（\path{dyntickRCU-irq-ssl.spin}）
以略多于五十万个状态产生了正确的验证，无错误通过。
然而，此版本的模型不处理 NMI，
这将在下一节中讨论。

\subsubsection{验证 NMI 处理程序}
\label{sec:formal:Validating NMI Handlers}

我们对 NMI 采取与中断相同的一般方法，
但要记住 NMI 不会嵌套。
这导致了如下所示的 \co{dyntick_nmi()} 进程：

\input{CodeSamples/formal/promela/dyntick/dyntickRCU-irq-nmi-ssl@dyntick_nmi.fcv}

当然，NMI 的存在要求对其他组件进行调整。
例如，\co{EXECUTE_MAINLINE()} 宏现在需要通过检查
\co{dyntick_nmi_done} 变量来同时关注 NMI 处理程序
（\co{in_dyntick_nmi}）和中断处理程序
（\co{in_dyntick_irq}），如下所示：

\input{CodeSamples/formal/promela/dyntick/dyntickRCU-irq-nmi-ssl@EXECUTE_MAINLINE.fcv}

我们还需要引入一个 \co{EXECUTE_IRQ()} 宏，
该宏检查 \co{in_dyntick_nmi} 以允许
\co{dyntick_irq()} 排除 \co{dyntick_nmi()}：

\input{CodeSamples/formal/promela/dyntick/dyntickRCU-irq-nmi-ssl@EXECUTE_IRQ.fcv}

还需要将 \co{dyntick_irq()} 转换为
\co{EXECUTE_IRQ()}，如下所示：

\input{CodeSamples/formal/promela/dyntick/dyntickRCU-irq-nmi-ssl@dyntick_irq.fcv}

\begin{fcvref}[ln:formal:promela:dyntick:dyntickRCU-irq-nmi-ssl:dyntick_irq]
请注意，我们已将 \qco{if} 语句直接展开编写
（例如，\clnrefrange{stmt1:b}{stmt1:e}）。
此外，处理严格局部状态的语句
（如 \clnref{inc_i}）不需要排除 \co{dyntick_nmi()}。
\end{fcvref}

最后，\co{grace_period()} 只需要少量更改：

\input{CodeSamples/formal/promela/dyntick/dyntickRCU-irq-nmi-ssl@grace_period.fcv}

\begin{fcvref}[ln:formal:promela:dyntick:dyntickRCU-irq-nmi-ssl:grace_period]
我们在 \clnref{MDL_NMI} 上为新的 \co{MAX_DYNTICK_LOOP_NMI}
参数添加了 \co{printf()}，并在 \clnref{nmi_done1,nmi_done2}
上将 \co{dyntick_nmi_done} 添加到 \co{shouldexit} 赋值中。
\end{fcvref}

此模型（\path{dyntickRCU-irq-nmi-ssl.spin}）
以数亿个状态产生了正确的验证，无错误通过。

\QuickQuiz{
	Paul \emph{总是}以这种痛苦的增量方式编写代码吗？
}\QuickQuizAnswer{
	不总是，但越来越频繁。
	在这个案例中，Paul 从包含中断处理程序的最小代码片段开始，
	因为他不确定如何最好地在 Promela 中建模中断。
	一旦让它运行起来，他就添加了其他特性。
	（但如果再做一次，他会从一个``玩具''处理程序开始。
	例如，他可能让处理程序将一个变量递增两次，
	并让主线代码验证该值始终为偶数。）

	为什么要用增量方法？
	请考虑以下这段话，出自Brian W. Kernighan：

	\begin{quote}
		调试的难度是编写代码的两倍。
		因此，如果你尽可能巧妙地编写代码，
		那么按定义，你就不够聪明来调试它。
	\end{quote}

	这意味着任何优化代码生产的尝试都应该将至少66\pct\ 的
	重点放在优化调试过程上，即使这会增加编码所花费的时间和
	精力。增量编码和测试是优化调试过程的一种方式，代价是
	编码工作量有所增加。
	Paul使用这种方法是因为他很少有奢侈的条件可以将整天
	（更不用说整周）用于编码和调试。
}\QuickQuizEnd

\subsubsection{（重新）学到的经验}
\label{sec:formal:Lessons (Re)Learned}

这次工作提供了一些（重新）学到的经验教训：

\begin{enumerate}
\item	{\bf Promela 和 Spin 可以验证中断/NMI 处理程序的交互。}
\item	{\bf 为代码编写文档可以帮助定位缺陷。}
	在本例中，文档编写工作定位了
	\co{rcu_enter_nohz()} 和 \co{rcu_exit_nohz()} 中
	一个放错位置的内存屏障，
	如以下补丁所示~\cite{PaulEMcKenney2008commit:ae66be9b71b1}。

\begin{VerbatimU}
 static inline void rcu_enter_nohz(void)
 {
+       mb();
        __get_cpu_var(dynticks_progress_counter)++;
-       mb();
 }

 static inline void rcu_exit_nohz(void)
 {
-       mb();
        __get_cpu_var(dynticks_progress_counter)++;
+       mb();
 }
\end{VerbatimU}

\item	{\bf 尽早、经常、彻底地验证你的代码。}
	这次工作定位了 \co{rcu_try_flip_waitack_needed()} 中
	一个非常微妙的缺陷，该缺陷很难通过测试或调试来发现，
	如以下补丁所示~\cite{PaulEMcKenney2008commit:d7c0651390b6}。

\begin{VerbatimU}
-       if ((curr - snap) > 2 || (snap & 0x1) == 0)
+       if ((curr - snap) > 2 || (curr & 0x1) == 0)
\end{VerbatimU}

\item	{\bf 始终验证你的验证代码。}
	通常的做法是故意插入一个缺陷，
	并验证验证代码能够捕获它。
	当然，如果验证代码未能捕获这个缺陷，
	你可能还需要验证缺陷本身，依此类推，
	无限递归。
	然而，如果你发现自己处于这种境地，
	好好睡一觉可能是一种极其有效的调试技术。
	然后你就会看到，验证验证代码的明显技术是
	在被验证的代码中故意插入缺陷。
	如果验证未能发现它们，那么验证代码显然有问题。
\item	{\bf 使用原子指令可以简化验证。}
	不幸的是，使用 \co{cmpxchg} 原子指令也会减慢
	关键的 \IXacr{irq} 快速路径，因此在本例中不适用。
\item	{\bf 需要复杂的形式验证通常表明需要重新思考你的设计。}
\end{enumerate}

关于最后一点，事实证明 dynticks 问题有一个简单得多的解决方案，
将在下一节中介绍。

\subsubsection{简单性避免了形式验证}
\label{sec:formal:Simplicity Avoids Formal Verification}

可抢占 RCU 的 dynticks 接口的复杂性主要是由于
\IRQ 和 NMI 使用相同的代码路径和相同的状态变量。
这引出了为 \IRQ 和 NMI 提供单独代码路径和变量的想法，
正如分层 RCU~\cite{PaulEMcKenney2008HierarchicalRCU}
所做的那样，这是由
Manfred Spraul~\cite{ManfredSpraul2008StateMachineRCU}
间接建议的。
这项工作在 v2.6.29 开发周期中被合入主线内核~\cite{PaulEMcKenney2008commit:64db4cfff99c}。

\subsubsection{简化 Dynticks 接口的状态变量}
\label{sec:formal:State Variables for Simplified Dynticks Interface}

\Cref{lst:formal:Variables for Simple Dynticks Interface}
显示了新的每 CPU 状态变量。
这些变量被分组到结构体中，以允许多个独立的 RCU 实现
（例如，\co{rcu} 和 \co{rcu_bh}）方便且高效地共享 dynticks 状态。
在以下内容中，它们可以被视为独立的每 CPU 变量。

\begin{listing}
\begin{VerbatimL}
struct rcu_dynticks {
	int dynticks_nesting;
	int dynticks;
	int dynticks_nmi;
};

struct rcu_data {
	...
	int dynticks_snap;
	int dynticks_nmi_snap;
	...
};
\end{VerbatimL}
\caption{简单 Dynticks 接口的变量}
\label{lst:formal:Variables for Simple Dynticks Interface}
\end{listing}

\co{dynticks_nesting}、\co{dynticks} 和 \co{dynticks_snap} 变量
用于 \IRQ\ 代码路径，\co{dynticks_nmi} 和
\co{dynticks_nmi_snap} 变量用于 NMI 代码路径，
尽管 NMI 代码路径也会引用（但不修改）
\co{dynticks_nesting} 变量。
这些变量的使用方式如下：

\begin{description}[style=nextline]
\item	[\tco{dynticks_nesting}]
	这个变量计数对应 CPU 应被监视 RCU 读端临界区的原因数量。
	如果 CPU 处于 dynticks-idle 模式，则它计数
	\IRQ\ 嵌套级别，否则比 \IRQ\ 嵌套级别大一。
\item	[\tco{dynticks}]
	如果对应 CPU 处于 dynticks-idle 模式且该 CPU 上
	没有正在运行的 \IRQ\ 处理程序，则该计数器的值为偶数，
	否则为奇数。
	换句话说，如果该计数器的值为奇数，
	则对应 CPU 可能处于 RCU 读端临界区中。
\item	[\tco{dynticks_nmi}]
	如果对应 CPU 处于 NMI 处理程序中，则该计数器的值为奇数，
	但仅当 NMI 到达时该 CPU 处于 dyntick-idle 模式
	且没有 \IRQ\ 处理程序在运行时才如此。
	否则，该计数器的值将为偶数。
\item	[\tco{dynticks_snap}]
	这将是 \co{dynticks} 计数器的快照，
	但仅当当前 RCU 宽限期已经延续了过长的时间时才如此。
\item	[\tco{dynticks_nmi_snap}]
	这将是 \co{dynticks_nmi} 计数器的快照，
	但同样仅当当前 RCU 宽限期已经延续了过长的时间时才如此。
\end{description}

如果 \co{dynticks} 和 \co{dynticks_nmi} 在给定时间间隔内
都取过偶数值，则对应 CPU 在该间隔内经历了一个\IX{quiescent state}。

\QuickQuiz{
	但如果一个 NMI 处理程序在一个 \IRQ\ 处理程序完成之前
	开始运行，并且该 NMI 处理程序一直运行到第二个
	\IRQ\ 处理程序开始，会怎样？
}\QuickQuizAnswer{
	这在单个 CPU 的范围内不可能发生。
	第一个 \IRQ\ 处理程序在 NMI 处理程序返回之前无法完成。
	因此，如果 \co{dynticks} 和 \co{dynticks_nmi}
	变量在给定时间间隔内都取过偶数值，
	则对应 CPU 确实在该间隔内的某个时刻处于静止状态。
}\QuickQuizEnd

\subsubsection{进入和离开 Dynticks-Idle 模式}
\label{sec:formal:Entering and Leaving Dynticks-Idle Mode}

\Cref{lst:formal:Entering and Exiting Dynticks-Idle Mode}
显示了 \co{rcu_enter_nohz()} 和 \co{rcu_exit_nohz()}，
它们进入和退出 dynticks-idle 模式，也称为``nohz''模式。
这两个函数从进程上下文中调用。

\begin{listing}
\begin{fcvlabel}[ln:formal:Entering and Exiting Dynticks-Idle Mode]
\begin{VerbatimL}[commandchars=\\\[\]]
void rcu_enter_nohz(void)
{
	unsigned long flags;
	struct rcu_dynticks *rdtp;

	smp_mb();			\lnlbl[mb]
	local_irq_save(flags);		\lnlbl[irq_sv]
	rdtp = &__get_cpu_var(rcu_dynticks); \lnlbl[get_ptr]
	rdtp->dynticks++;		\lnlbl[inc_cnt]
	rdtp->dynticks_nesting--;	\lnlbl[dec_nst]
	WARN_ON(rdtp->dynticks & 0x1);
	local_irq_restore(flags);	\lnlbl[irq_rs]
}

void rcu_exit_nohz(void)
{
	unsigned long flags;
	struct rcu_dynticks *rdtp;

	local_irq_save(flags);
	rdtp = &__get_cpu_var(rcu_dynticks);
	rdtp->dynticks++;
	rdtp->dynticks_nesting++;
	WARN_ON(!(rdtp->dynticks & 0x1));
	local_irq_restore(flags);
	smp_mb();
}
\end{VerbatimL}
\end{fcvlabel}
\caption{进入和退出 Dynticks-Idle 模式}
\label{lst:formal:Entering and Exiting Dynticks-Idle Mode}
\end{listing}

\begin{fcvref}[ln:formal:Entering and Exiting Dynticks-Idle Mode]
\Clnref{mb} 确保任何先前的内存访问（可能包括来自 RCU
读端临界区的访问）在标记进入 dynticks-idle 模式之前被其他 CPU 看到。
\Clnref{irq_sv,irq_rs} 禁用和重新启用 \IRQ。
\Clnref{get_ptr} 获取指向当前 CPU 的 \co{rcu_dynticks}
结构的指针，
\clnref{inc_cnt} 递增当前 CPU 的 \co{dynticks} 计数器，
鉴于我们在进程上下文中进入 dynticks-idle 模式，
该计数器现在应该是偶数。
最后，\clnref{dec_nst} 递减 \co{dynticks_nesting}，
现在应该为零。
\end{fcvref}

\co{rcu_exit_nohz()} 函数非常相似，但递增 \co{dynticks_nesting}
而非递减，并检查相反的 \co{dynticks} 极性。

\subsubsection{来自 Dynticks-Idle 模式的 NMI}
\label{sec:formal:NMIs From Dynticks-Idle Mode}

\begin{fcvref}[ln:formal:NMIs From Dynticks-Idle Mode]
\Cref{lst:formal:NMIs From Dynticks-Idle Mode}
显示了 \co{rcu_nmi_enter()} 和 \co{rcu_nmi_exit()} 函数，
它们分别通知 RCU 从 dynticks-idle 模式进入和退出 NMI。
然而，如果 NMI 在 \IRQ\ 处理程序期间到达，
那么 RCU 已经在监视来自该 CPU 的 RCU 读端临界区，
因此如果 \co{dynticks} 为奇数，\co{rcu_nmi_enter()} 的
\clnref{chk_ext1,ret1} 和 \co{rcu_nmi_exit()} 的
\clnref{chk_ext2,ret2} 会静默返回。
否则，两个函数递增 \co{dynticks_nmi}，
\co{rcu_nmi_enter()} 使其保持奇数值，
\co{rcu_nmi_exit()} 使其保持偶数值。
两个函数分别在 \clnref{mb1,mb2} 上在此递增和可能的 RCU
读端临界区之间执行内存屏障。
\end{fcvref}

\begin{listing}
\begin{fcvlabel}[ln:formal:NMIs From Dynticks-Idle Mode]
\begin{VerbatimL}[commandchars=\\\[\]]
void rcu_nmi_enter(void)
{
	struct rcu_dynticks *rdtp;

	rdtp = &__get_cpu_var(rcu_dynticks);
	if (rdtp->dynticks & 0x1)	\lnlbl[chk_ext1]
		return;			\lnlbl[ret1]
	rdtp->dynticks_nmi++;
	WARN_ON(!(rdtp->dynticks_nmi & 0x1));
	smp_mb();			\lnlbl[mb1]
}

void rcu_nmi_exit(void)
{
	struct rcu_dynticks *rdtp;

	rdtp = &__get_cpu_var(rcu_dynticks);
	if (rdtp->dynticks & 0x1)	\lnlbl[chk_ext2]
		return;			\lnlbl[ret2]
	smp_mb();			\lnlbl[mb2]
	rdtp->dynticks_nmi++;
	WARN_ON(rdtp->dynticks_nmi & 0x1);
}
\end{VerbatimL}
\end{fcvlabel}
\caption{来自 Dynticks-Idle 模式的 NMI}
\label{lst:formal:NMIs From Dynticks-Idle Mode}
\end{listing}

\subsubsection{来自 Dynticks-Idle 模式的中断}
\label{sec:formal:Interrupts From Dynticks-Idle Mode}

\begin{fcvref}[ln:formal:Interrupts From Dynticks-Idle Mode]
\Cref{lst:formal:Interrupts From Dynticks-Idle Mode}
显示了 \co{rcu_irq_enter()} 和 \co{rcu_irq_exit()}，
它们分别通知 RCU 进入和退出 \IRQ\ 上下文。
\co{rcu_irq_enter()} 的 \clnref{inc_nst} 递增 \co{dynticks_nesting}，
如果该变量已经非零，\clnref{ret1} 会静默返回。
否则，\clnref{inc_dynt1} 递增 \co{dynticks}，
该计数器随后将具有奇数值，与该 CPU 现在可以执行
RCU 读端临界区的事实一致。
因此 \clnref{mb1} 执行一个内存屏障以确保
\co{dynticks} 的递增在后续 \IRQ\ 处理程序可能执行的
任何 RCU 读端临界区之前被看到。

\begin{listing}
\begin{fcvlabel}[ln:formal:Interrupts From Dynticks-Idle Mode]
\begin{VerbatimL}[commandchars=\\\[\]]
void rcu_irq_enter(void)
{
	struct rcu_dynticks *rdtp;

	rdtp = &__get_cpu_var(rcu_dynticks);
	if (rdtp->dynticks_nesting++)		\lnlbl[inc_nst]
		return;				\lnlbl[ret1]
	rdtp->dynticks++;			\lnlbl[inc_dynt1]
	WARN_ON(!(rdtp->dynticks & 0x1));
	smp_mb();				\lnlbl[mb1]
}

void rcu_irq_exit(void)
{
	struct rcu_dynticks *rdtp;

	rdtp = &__get_cpu_var(rcu_dynticks);
	if (--rdtp->dynticks_nesting)		\lnlbl[dec_nst]
		return;				\lnlbl[ret2]
	smp_mb();				\lnlbl[mb2]
	rdtp->dynticks++;			\lnlbl[inc_dynt2]
	WARN_ON(rdtp->dynticks & 0x1);		\lnlbl[chk_even]
	if (__get_cpu_var(rcu_data).nxtlist ||	\lnlbl[chk_cb:b]
	    __get_cpu_var(rcu_bh_data).nxtlist)
		set_need_resched();		\lnlbl[chk_cb:e]
}
\end{VerbatimL}
\end{fcvlabel}
\caption{来自 Dynticks-Idle 模式的中断}
\label{lst:formal:Interrupts From Dynticks-Idle Mode}
\end{listing}

\co{rcu_irq_exit()} 的 \clnref{dec_nst} 递减 \co{dynticks_nesting}，
如果结果非零，\clnref{ret2} 静默返回。
否则，\clnref{mb2} 执行内存屏障以确保
\clnref{inc_dynt2} 上 \co{dynticks} 的递增在先前
\IRQ\ 处理程序可能执行的任何 RCU 读端临界区之后被看到。
\Clnref{chk_even} 验证 \co{dynticks} 现在为偶数，
与 dynticks-idle 模式中不能出现 RCU 读端临界区的事实一致。
\Clnrefrange{chk_cb:b}{chk_cb:e} 检查先前的 \IRQ\ 处理程序
是否排入了任何 RCU 回调，如果是则通过重调度 API
强制此 CPU 退出 dynticks-idle 模式。
\end{fcvref}

\subsubsection{检查 Dynticks 静止状态}
\label{sec:formal:Checking For Dynticks Quiescent States}

\begin{fcvref}[ln:formal:Saving Dyntick Progress Counters]
\Cref{lst:formal:Saving Dyntick Progress Counters}
显示了 \co{dyntick_save_progress_counter()}，它获取
指定 CPU 的 \co{dynticks} 和 \co{dynticks_nmi} 计数器的快照。
\Clnref{snap,snapn} 将这两个变量快照到局部变量中，
\clnref{mb} 执行内存屏障以与
\cref{lst:formal:Entering and Exiting Dynticks-Idle Mode,%
lst:formal:NMIs From Dynticks-Idle Mode,%
lst:formal:Interrupts From Dynticks-Idle Mode}
中函数的内存屏障配对。
\Clnref{rec_snap,rec_snapn} 记录快照以供后续调用
\co{rcu_implicit_dynticks_qs()} 使用，
\clnref{chk_prgs} 检查 CPU 是否处于 dynticks-idle 模式且
没有正在进行的 \IRQ 或 NMI（换句话说，两个快照都是偶数值），
因此处于扩展静止状态。
如果是，\clnref{cnt:b,cnt:e} 计数此事件，
\clnref{ret} 在 CPU 处于静止状态时返回 true。
\end{fcvref}

\begin{listing}
\begin{fcvlabel}[ln:formal:Saving Dyntick Progress Counters]
\begin{VerbatimL}[commandchars=\\\[\]]
static int
dyntick_save_progress_counter(struct rcu_data *rdp)
{
	int ret;
	int snap;
	int snap_nmi;

	snap = rdp->dynticks->dynticks;		\lnlbl[snap]
	snap_nmi = rdp->dynticks->dynticks_nmi;	\lnlbl[snapn]
	smp_mb();				\lnlbl[mb]
	rdp->dynticks_snap = snap;		\lnlbl[rec_snap]
	rdp->dynticks_nmi_snap = snap_nmi;	\lnlbl[rec_snapn]
	ret = ((snap & 0x1) == 0) && ((snap_nmi & 0x1) == 0); \lnlbl[chk_prgs]
	if (ret)				\lnlbl[cnt:b]
		rdp->dynticks_fqs++;		\lnlbl[cnt:e]
	return ret;				\lnlbl[ret]
}
\end{VerbatimL}
\end{fcvlabel}
\caption{保存 Dyntick 进度计数器}
\label{lst:formal:Saving Dyntick Progress Counters}
\end{listing}

\begin{fcvref}[ln:formal:Checking Dyntick Progress Counters]
\Cref{lst:formal:Checking Dyntick Progress Counters}
显示了 \co{rcu_implicit_dynticks_qs()}，它被调用以检查
CPU 是否在调用 \co{dynticks_save_progress_counter()} 之后
进入了 dyntick-idle 模式。
\Clnref{curr,currn} 获取对应 CPU 的 \co{dynticks} 和
\co{dynticks_nmi} 变量的新快照，
而 \clnref{snap,snapn} 检索之前由
\co{dynticks_save_progress_counter()} 保存的快照。
\Clnref{mb} 然后执行内存屏障以与
\cref{lst:formal:Entering and Exiting Dynticks-Idle Mode,%
lst:formal:NMIs From Dynticks-Idle Mode,%
lst:formal:Interrupts From Dynticks-Idle Mode}
中函数的内存屏障配对。
\Clnrefrange{chk_q:b}{chk_q:e} 然后检查 CPU 是否当前
处于静止状态（\co{curr} 和 \co{curr_nmi} 具有偶数值），
或者自上次调用 \co{dynticks_save_progress_counter()}
以来是否经历了静止状态（\co{dynticks} 和 \co{dynticks_nmi}
的值已改变）。
如果这些检查确认 CPU 已经经历了 dyntick-idle 静止状态，
则 \clnref{cnt} 计数该事实，
\clnref{ret_1} 返回该事实的指示。
无论如何，\clnref{chk_race} 检查可能导致 RCU
等待已离线 CPU 的竞态条件。
\end{fcvref}

\begin{listing}
\begin{fcvlabel}[ln:formal:Checking Dyntick Progress Counters]
\begin{VerbatimL}[commandchars=\\\[\]]
static int
rcu_implicit_dynticks_qs(struct rcu_data *rdp)
{
	long curr;
	long curr_nmi;
	long snap;
	long snap_nmi;

	curr = rdp->dynticks->dynticks;		\lnlbl[curr]
	snap = rdp->dynticks_snap;		\lnlbl[snap]
	curr_nmi = rdp->dynticks->dynticks_nmi;	\lnlbl[currn]
	snap_nmi = rdp->dynticks_nmi_snap;	\lnlbl[snapn]
	smp_mb();				\lnlbl[mb]
	if ((curr != snap || (curr & 0x1) == 0) && \lnlbl[chk_q:b]
	    (curr_nmi != snap_nmi || (curr_nmi & 0x1) == 0)) { \lnlbl[chk_q:e]
		rdp->dynticks_fqs++;		\lnlbl[cnt]
		return 1;			\lnlbl[ret_1]
	}
	return rcu_implicit_offline_qs(rdp);	\lnlbl[chk_race]
}
\end{VerbatimL}
\end{fcvlabel}
\caption{检查 Dyntick 进度计数器}
\label{lst:formal:Checking Dyntick Progress Counters}
\end{listing}

\QuickQuiz{
	这仍然相当复杂。
	为什么不直接使用一个带有每 CPU 位的 \co{cpumask_t}，
	在进入 \IRQ\ 或 NMI 处理程序时清除该位，
	在退出时设置它？
}\QuickQuizAnswer{
	虽然这种方法在功能上是正确的，但它会导致
	在大型机器上产生过多的 \IRQ\ 进入/退出开销。
	相比之下，本节中阐述的方法允许每个 CPU 在
	\IRQ\ 和 NMI 进入/退出时只触碰每 CPU 数据，
	从而产生更低的 \IRQ\ 进入/退出开销，
	尤其是在大型机器上。
}\QuickQuizEnd

Linux 内核 RCU 的 dyntick-idle 代码后来又基于
Andy Lutomirski~\cite{PaulMcKenney2015dyntickAndyNMI}
的建议被重写了一次，
但现在是时候总结并继续讨论其他主题了。

\subsubsection{讨论}
\label{sec:formal:Discussion}

一个小小的观点转变导致了 RCU dynticks 接口的实质性简化。
导致这种简化的关键变化是将 \IRQ\ 和 NMI 上下文之间的共享减到最小。
在这个简化接口中，唯一的共享是 NMI 上下文对
\IRQ\ 变量（\co{dynticks} 变量）的引用。
这种类型的共享是良性的，因为 NMI 函数从不更新该变量，
因此其值在 NMI 处理程序的整个生命周期中保持不变。
这种共享的限制使得各个函数可以逐个理解，
与\cref{sec:formal:Promela Parable: dynticks and Preemptible RCU}
中描述的情况形成鲜明对比——在那里，NMI 可能在
\IRQ\ 函数执行期间的任何时刻修改共享状态。

验证是一件非常好的事，但简单性更好。

% Restore frame around VerbatimN in two-column layout
\IfTwoColumn{
\RecustomVerbatimEnvironment{VerbatimN}{Verbatim}%
{numbers=left,numbersep=3pt,xleftmargin=5pt,xrightmargin=5pt,frame=single}
\setlength{\lnlblraise}{6pt}
}{}
